% ======================================================================
%  PAPER I of the two-paper split
%  Controlled-rail brachistochrones in non-stationary spacetimes:
%  conformal symmetry, Kodama energy, and Vaidya dynamics
%  IOP iopart class -- Classical and Quantum Gravity
%  (spherically symmetric sector: FLRW + Vaidya; shared control-theory
%   foundations live here and are cited by Paper II)
% ======================================================================
\documentclass[12pt]{iopart}

% iopart pre-defines \equation*, which clashes with amsmath; free it first
\makeatletter
\expandafter\let\csname equation*\endcsname\relax
\expandafter\let\csname endequation*\endcsname\relax
\makeatother
\usepackage{amsmath}
\usepackage{iopams}
\usepackage{graphicx}
\usepackage{bm}
\usepackage{xcolor}
\usepackage{hyperref}
\hypersetup{colorlinks=true, linkcolor=blue, citecolor=blue, urlcolor=blue}

% figures shared with the master/Paper II
\graphicspath{{../paper/Immagini/}}

% prose placeholder helper
\newcommand{\todo}[1]{{\color{red}\small[[#1]]}}

% shortcuts
\newcommand{\Ehat}{\hat E}
\newcommand{\rp}{r_{+}}
\newcommand{\re}{r_{e}}
\newcommand{\Jc}{J_{c}}
\newcommand{\Af}{A_{\mathrm{freeze}}}
\newcommand{\dd}{\mathrm{d}}

\begin{document}

\title{Controlled-rail brachistochrones in non-stationary spacetimes:
conformal symmetry, Kodama energy, and Vaidya dynamics}

\author{Iman Rosignoli}
\address{University of Pavia, Pavia, Italy}
\ead{iman.rosignoli01@universitadipavia.it}

\begin{abstract}
\todo{Paper I abstract (<=300 words, CQG): rail invariant -u.W=Ehat via
Pontryagin; W-hierarchy Killing -> CKV -> Kodama; FLRW degenerate base and
Vaidya dynamical black hole (Kodama energy conserved on the rail,
endpoint-sensitive costate memory, bounce regularization); Vaidya first-order
adiabatic correction and the mass-function universal source S\_D=[r p\_r]-lambda.
Status: existence/normality/HJB-verification; outgoing-Vaidya BVP left open.}
\end{abstract}

\maketitle
{\hypersetup{linkcolor=black}\tableofcontents}

% ======================================================================
\section{Introduction}
\todo{Adapt master intro; scope to the spherically symmetric sector; forward
Paper II for the rotating (Thakurta-Kerr) case.}

% ---- SHARED FOUNDATIONS (cited by Paper II) --------------------------
\section{The controlled-rail principle}
\todo{Migrate master Sec. 'The controlled-rail principle':
2.1 rail constraint + indicatrix + optimal-control problem statement;
2.2 existence, normality, non-autonomous HJB verification;
2.3 W-hierarchy Killing--CKV--Kodama.}

\section{FLRW: the degenerate base case}
\todo{Migrate master FLRW section.}

\section{Vaidya: dynamical black hole and the Kodama rail}
\todo{Migrate master Vaidya section: Kodama energy invariant; phenomenology
(penetration, timing, bounce); plunge-radius law and the (absent) evaporative
inversion with the negative-rate-continuation scoping; closed-form adiabatic
correction and horizon dilogarithm; the mass-function universal source
S\_D=[r p\_r]-lambda.}

\section{Conclusions}
\todo{Three blocks: Proved / Conditional and numerical / Outlook (Vaidya scope,
outgoing-BVP open).}

\appendix
\section{Validation and consistency checks}
\todo{Vaidya/shared parts of the master validation appendix.}

\section{Explicit adiabatic forms on the genus-degeneration loci (Vaidya)}
\todo{Vaidya J\_deg (algebraic degeneration, r\_d<0) only; NOT called a physical
separatrix.}

\section{Complete first-order extended-Hamiltonian correction (Vaidya)}
\todo{Canonical derivation in (r,p\_r); mass-function source.}

% ---- back matter -----------------------------------------------------
\section*{Data availability}
\todo{Same Zenodo DOI 10.5281/zenodo.21707378 statement as the master.}

\ack
\todo{Same acknowledgements as master.}

\nocite{*}
\bibliographystyle{iopart-num}
\bibliography{../paper/refs}

\end{document}
